\documentclass[11pt]{article}

\usepackage[utf8]{inputenc}
\usepackage[T1]{fontenc}
\usepackage{lmodern}
\usepackage[french]{babel}   
\usepackage{float}
\usepackage[skip=0.5cm]{parskip}
\usepackage{amsmath,amssymb,amsfonts}

\author{Jean BAYLON}
\title{Une grille d'analyse générique : entités, structures, et causalité}
\date{Version du \today{}}

\begin{document}
	\maketitle
	\abstract{La mise en place d'un modèle est une tâche complexe : peu de concepts, indépendants, avec un pouvoir explicatif ou prédictif élevé. La description de méta modèles est, à l'évidence, d'autant plus exigeante qu'elle exige une forme d'exhaustivité. Nous proposons une formalisation se basant sur trois notions : \textbf{les entités s'organisent en structures qui évoluent par causalité}. Notre travail est issu d'une recherche personnelle. Nous le mettons en ligne dans l'espoir d'être utile.}
	
	\section{Le modèle général}
	
		Dans sa forme la plus générale, notre modèle postule simplement qu'il existe des entités qui évoluent par une logique de causalité pour un observateur. Si la mesure précise du temps n'est pas forcément nécessaire, il est par contre indispensable de savoir situer l'avant de l'après. L'introduction de la causalité répond à cette problématique. Formellement, elle ordonne la succession, pour un observateur, des observations suivant une logique implacable : personne ne peut constater l'effet avant la cause. 
		
		\begin{table}[H]
			\centering
			\begin{tabular}{|l|l|}
				\hline 
				Nom & Description \\
				\hline 
				Entité & Ce qui est observable, caractérisé par une chaine d'événements \\
				\hline 
				Structure & Le réseau invariant (pour les entités) et persistent \\
				\hline 
				Causalité & La dynamique \\
				\hline 
			\end{tabular}
			\caption{Résumé du modèle}
		\end{table}
		
		
		\subsection{États et interactions : le concept clé d'événement}
		
			Le concept de causalité permet de déduire deux nouvelles notions : 
			\begin{itemize}
				\item dans sa dimension statique, la causalité caractérise l'état, la configuration 
				\item dans sa dimension dynamique, elle décrit comment des interactions modifient le système
			\end{itemize}
			
			
			\vspace{0.5cm}
			
			Finalement, le lien entre entité et interaction est fort: 
			\begin{enumerate}
				\item une entité est transformée par les interactions qu'elle subit
				\item la présence de l'entité ou sa transformation peut émettre en retour des interactions
			\end{enumerate}
		
		
			Cette vue qu'on pourrait qualifier d'immédiate, de \og{}bon sens\fg{} est souvent suffisante, et nous en faisons l'expérience immédiate. On sait décrire telle ou telle caractéristique de la chose observée. Une telle approche, cependant, connait des limites fortes, très tôt dans son usage : 
			\begin{enumerate}
				\item elle ne résiste pas au temps, au changement qu'il imprime sur l'entité. Ce n'est pas l'état qui nous permet de reconnaitre une entité, mais notre capacité à l'identifier au delà de ses caractéristiques immédiates. Autrement dit, l'état passe, mais c'est l'identité qui reste. 
				\item un changement de comportement intrinsèque, une restructuration interne, sont par essence des changements propres au système étudié. De fait, ils ne sauraient constituer une interaction en l'absence de déclencheur externe. La notion clé est l'événement, dont les interactions sont une facette. 
			\end{enumerate}
			
			Ainsi, fondamentalement, ce qui caractérise une entité, ou un système, ou tout observable, est sa chaine causale : 
			\begin{enumerate}
				\item l'identité a été acquise à sa création, ce qui constitue son événement fondateur
				\item ses états s'enchainent, ce qui constitue son historique causal. 
				\item la fin de l'entité est un événement également
			\end{enumerate}
			
			En traitant d'abord les événements, les entités sont des observables temporaires. Qu'ils soient individuellement présents ou non ne marque ni le début ni la fin de l'étude ou de l'observation. \textbf{Une entité a alors une identité propre quand elle reste une chaine causale d'événements ininterrompue : on sait l'expliquer et la reconnaitre}. 
		
		\subsection{Le temps n'est pas un concept premier}
		
			Le temps n'est plus un concept premier, il donne une solution numérique, mesurable, à des phénomènes d'observations coordonnées (avant / après / simultanément). Ce changement de paradigme permet d'adresser la problématique des temps différents en fonction de l'échelle d'un système observé. Par exemple, notre corps change plus vite que nous, mais nous conservons notre identité. Une société a un temps plus long que ses membres, car un changement d'état implique la synchronisation de nombreux éléments (opportunités sociales, faiblesse de la structure existante, volonté interne ou externe de changement). On peut donc considérer une unité de temps fondamentale, une sorte de \textit{tick} comme un changement perceptible dans l'état du système observé. Ce tick peut être, chez une personne, un mariage, un nouvel emploi, une étape clé de sa vie. Au niveau d'une société, ce peut être une nouvelle constitution, une rupture de politique générale ayant des effets observables. Mais dans tous les cas, le temps n'est pas uniforme, le temps n'est pas un concept primordial. Bien évidemment, les concepts sont très fortement intriqués : parler de transition implique un temps, ce qui amène à des débats très pointus que nous esquivons ici. 
		
		\subsection{Les relations causales ne sont pas toutes déterministes}
		
			Détaillons à présent la logique de passage d'un état à un autre, c'est à dire la transition par les interactions. La causalité est donc une relation d'accessibilité : peut on passer de cet état à celui-ci ? 
			Cependant, le mécanisme n'a rien d'automatique ou de déterministe, on choisira la relation la plus adaptée parmi : 
			\begin{itemize}
				\item une causalité directe, qu'on retrouve dans les systèmes physiques à notre échelle. Sa forme la plus canonique est par exemple $\vec{F} = m\vec{a}$. 
				\item une causalité probabiliste, qu'on retrouve en mécanique quantique. On sait que l'état suivant sera pris parmi ceux là, avec une distribution de probabilité. 
				\item une causalité tendancielle, qu'on retrouve plus souvent dans les sciences sociales, comme des comportements préférentiellement choisis. Par exemple, le mécanisme d'analogie historique revient souvent en diplomatie. 
			\end{itemize}
		
		\subsection{Le rôle clé de l'observateur}
		
			Fondamentalement, il s'agit de répondre à la question de comment un système, dans un état initial donné, évolue suivant une relation de causalité plus ou moins complexe. Il est impératif de se rappeler que l'observateur fait partie du modèle, c'est lui qui caractérise ces états et qui cherche à les lier entre eux. Il n'y a donc pas d'état initial, au sens d'une réalité ontologique de l'univers. C'est bien une coupure arbitraire décidée par l'observateur. En pratique, c'est l'analyste du système qui décrète qu'il commence son observation en ce point là d'espace temps. \textbf{L'état initial est donc toujours un état hérité, que l'on choisit de figer pour commencer le calcul ou l'étude.} 
			
		
	
	\section{Distinguer le tout de la partie}
	
		Ce que l'observateur veut comprendre n'est pas nécessairement ce qu'il peut observer directement. Par exemple, un sociologue décrit une société en observant un nombre réduit de ses membres. Un physicien spécialiste de physique quantique se retrouve à utiliser un bagage théorique notoirement contre-intuitif à notre échelle. Pour adresser ces questions, le paradigme élément - système ou entité - structure s'est montré efficace. Cette section en détaille certains aspects. Mais remarquons d'abord qu'on peut tout à fait l'appliquer plusieurs fois, pour avoir une hiérarchie à différents niveaux, comme par exemple monde, société, communautés, individus.
	
					
		\subsection{La constitution de la structure : un réseau avant tout}
		
			La structure porte en elle une information clé : elle reste organisée quand ses entités changent une à une. Très simplement, une entreprise se maintient quand ses employés la quittent, notre corps se renouvelle mais nous restons nous même. C'est donc, au delà de chaque entité spécifique, le réseau des relations qu'elles entretiennent qui nous éclaire sur ce qu'est la structure, sur comment elle marche. Ce sont les propriétés des éléments, les natures des relations qu'ils tissent, qui constituent l'essence de la structure. Nous invitons tout utilisateur du modèle à se rappeler que nous n'avons pas imposé de nature spécifique à notre notion d'entité : ce peut être de la matière, de l'information, une relation, etc.
		
		\subsection{Les propriétés de la structure : un cadre non neutre}
		
			Considérer la structure comme un simple support physique ou un cadre vide pour les éléments est trop réducteur. Souvent, elle s'adapte dynamiquement à ce qu'elle contient et influence activement les éléments :
			\begin{enumerate}
				\item \textbf{L'émergence :} la structure est un sujet d'étude en soi. L'émergence d'un comportement global peut ne pas se réduire aux actions de ses éléments (ex: la formation de comportements intelligents dans les réseaux de neurones).
				\item \textbf{La co-évolution :} il existe des cas où les propriétés des éléments modifient la structure. L'exemple typique est physique : la matière courbe l'espace-temps, qui, en retour, oriente les déplacements de la matière.
				\item \textbf{Le filtre actif :} la structure n'est pas un support de transmission neutre. Sur les réseaux sociaux, par exemple, l'intermédiation algorithmique met en avant certains contenus et en rend d'autres invisibles. L'algorithme est précisément l'application d'entreprise et la source de monétisation de ces autoroutes de l'information.
			\end{enumerate}
			
		
			
	
	\section{Conclusion} 
			
		Bien évidemment, le modèle que nous proposons ici est extrêmement général. Si notre inspiration physique initiale est claire, nous pensons que ce vocabulaire est suffisamment abstrait pour être pertinent dans un cadre d'analyse sociologique ou systémique. Sur le plan pratique, l'application de cette grille de lecture requiert une méthode stricte, consistant à identifier :
		
		\begin{enumerate}
			\item \textbf{La finalité et le cadre de référence} : expliciter les objectifs du modèle et les biais de l'analyste qui pourraient impacter les conclusions. La coupure arbitraire de l'état initial gagne souvent à être enrichie d'une contextualisation permettant de saisir pourquoi ce moment serait clé dans l'étude proposée. 
			\item \textbf{La cartographie du système} : isoler les observables fondamentaux pour déterminer ce qui relève de l'entité, ce qui constitue la structure, et la nature des relations qui les unissent. 
			\item \textbf{La qualification des chaînes causales} : définir les critères permettant de distinguer deux états successifs et caractériser la nature des interactions (déterministes, probabilistes ou tendancielles). Cette analyse peut s'appuyer sur des outils statistiques permettant de mesurer, même dans la définition de tendances, des indicateurs dits \og{}proxy\fg{} et proposer une validation des éléments constatés (donc, après les avoir observés). 
		\end{enumerate}
		
				
		Une fois cette méthodologie posée, ce cadre formalisé permet de soulever des questions particulièrement prometteuses. Par exemple :
		\begin{itemize}
			\item Une analyse tendancielle des indicateurs économiques permettrait-elle d'anticiper des ruptures nettes dans la structure sociale, en fonction des interactions d'acteurs clés ?
			\item Quelles applications sont possibles dans l'étude des réseaux sociaux ? Comment la structure (l'algorithme) peut-elle être modélisée pour contrer l'ingérence, et peut-on détecter mathématiquement des changements topologiques caractéristiques de manipulations comme l'astroturfing ?
		\end{itemize}
		
	
\end{document}